\documentclass[11pt]{article}
\usepackage{lmodern}
\usepackage[T1]{fontenc}
\usepackage{amsmath,amssymb,amsthm}
\usepackage[margin=1in]{geometry}
\usepackage[colorlinks=true,urlcolor=blue,citecolor=blue,linkcolor=blue]{hyperref}

\newtheorem{theorem}{Theorem}[section]
\newtheorem{proposition}[theorem]{Proposition}
\newtheorem{lemma}[theorem]{Lemma}
\newtheorem{corollary}[theorem]{Corollary}
\newtheorem{remark}[theorem]{Remark}
\newcommand{\Q}{\mathbb{Q}}
\newcommand{\R}{\mathbb{R}}
\newcommand{\Z}{\mathbb{Z}}
\newcommand{\F}{\mathbb{F}}
\newcommand{\calV}{\mathcal{V}}
\newcommand{\calB}{\mathcal{B}}
\newcommand{\ord}{\operatorname{ord}}
\newcommand{\Reg}{\operatorname{Reg}}
\newcommand{\Res}{\operatorname*{Res}}
\newcommand{\vol}{\operatorname{vol}}

\title{Rational-Point Height Singularities and Quadratic-Character Series Attached to Elliptic Curves}
\author{\small \href{mailto:admin@cosmosgeek.org}{admin@cosmosgeek.org}\qquad \href{mailto:b321014@iiit-bh.ac.in}{b321014@iiit-bh.ac.in}}
\date{Working draft}

\begin{document}
\maketitle

\begin{abstract}
For a rational elliptic curve presented as a double cover of the projective line, we attach the positive squareclasses of a binary quartic, a rational-square height series, and a two-variable series of primitive quadratic Dirichlet $L$-functions.  We prove an exact identification of the rational-square locus with rational elliptic points, a regulator-normalized height-counting asymptotic, and its Abelian pole-type singularity.  Two specializations of the quadratic-character series recover the same height series: evaluation at character parameter zero and the residue at character parameter one.  Independently, we prove an Euler-product formula for the bulk average of all locally admissible slopes, and a meromorphic continuation of that \emph{bulk coefficient} to character parameter zero.  Its order is $1+b-a$, where $b$ is the number of rational Galois orbits of the four branch points and $a$ is the complex analytic rank of the elliptic curve.  We identify a rigorous obstruction to equating this order with the height singularity: all rational-square slopes may be deleted without changing the bulk coefficient, and the normalized finite-height means have a divergent first character derivative at zero.  The two asymptotics therefore do not prove Birch--Swinnerton-Dyer.  We also compute a self-dual Poisson sum, local squareclass laws in two integral models, and the signed Frobenius and height series on the symmetric square. No finiteness assumption on the Tate--Shafarevich group is made.
\end{abstract}

\section{Introduction and scope}
Let $E/\Q$ be an elliptic curve with a rational Weierstrass model
\begin{equation}\label{eq:model}
 E:\quad y^2=f(x),\qquad
 f(x)=A x^3+B x^2+C x+D\in\Z[x],\quad A\ne0,
\end{equation}
where $f$ is separable.  The point at infinity makes this genus-one curve an elliptic curve.  We use the binary quartic
\begin{equation}\label{eq:quartic}
 G(a,b)=b^4f(a/b)
       =b(Aa^3+Ba^2b+Cab^2+Db^3).
\end{equation}
The four geometric branch points of the degree-two map $x:E\to\mathbb{P}^1$ are the three roots of $f$ and infinity.

Let $r=\operatorname{rank}E(\Q)$ and
\[
 a_E=\ord_{s=1}L(E,s).
\]
The Birch--Swinnerton-Dyer rank conjecture asserts $r=a_E$.  The aim here is \emph{not} to establish that equality, but to give a precise, self-contained account of what two natural singularity calculations actually prove and where their comparison fails.

Fix a finite set $S$ of rational primes containing $2$, all primes dividing the leading coefficient or discriminant of $f$ and all primes of bad reduction of \eqref{eq:model}, and, for the uniform Euler-product theorem below, $3,5,7$.  The extra three small primes are a technical convenience ensuring that the remaining local correction factors are nonzero at the boundary.  One may omit them on an individual curve after checking the corresponding factors directly.

Let $\calV=\calV_{E,S}$ consist of primitive pairs
\begin{equation}\label{eq:slopes}
 (a,b)\in\Z\times\Z_{>0},\qquad \gcd(a,b)=1,\qquad
 G(a,b)>0,\qquad G(a,b)\in\Q_p^{\times2}\ (p\in S).
\end{equation}
Write $H(a,b)=(a^2+b^2)^{1/2}$.  There is a unique positive fundamental quadratic discriminant $D_{a,b}$ such that $G(a,b)/D_{a,b}$ is a square in $\Q^\times$; set $D_{a,b}=1$ precisely when $G(a,b)$ is a rational square.  Since $G(a,b)$ is an integer, the latter means an integer square.  Let $\chi_D$ be the corresponding primitive even quadratic Dirichlet character, with $L(w,\chi_1)=\zeta(w)$.

We use the modularity theorem \cite{BCDT} to regard $L(E,s)$ as the entire weight-two elliptic $L$-function, including its bad Euler factors.  We use only the Mordell--Weil theorem, the height comparison, classical quadratic Dirichlet functional equations, and elementary adelic averaging for the arithmetic assertions; see \cite{SilvermanHeight,SilvermanAEC,DirichletClass,HeckeTate}.  In particular no special-value formula of BSD type is an input.

\section{Rational square slopes and the height singularity}
\subsection{Exact parametrization}
\begin{lemma}\label{lem:points}
For $(a,b)\in\calV$, the condition $D_{a,b}=1$ is equivalent to the existence of a finite rational point $P\in E(\Q)$ with $x(P)=a/b$ and $y(P)\ne0$.  Each such rational $x$-coordinate is counted once by $\calV$ and corresponds to the two points $P,-P$.  All rational finite points with $y(P)\ne0$ occur, independently of the choice of $S$.
\end{lemma}
\begin{proof}
If $G(a,b)=c^2$ with $c\in\Z\setminus\{0\}$, then
$f(a/b)=(c/b^2)^2$, giving $(a/b,\pm c/b^2)\in E(\Q)$.
Conversely, if $x(P)=a/b$ in lowest terms, then
$G(a,b)=(b^2y(P))^2$.  Since $G(a,b)\in\Z$, its rational square root is integral.  A global rational square is a square in each $\Q_p$, so the local conditions in \eqref{eq:slopes} remove no such points.  The sign convention $b>0$ makes the primitive pair unique.  The points with $y=0$ and the point at infinity are finitely many and are omitted by definition.
\end{proof}

\subsection{Height normalization and exact leading constant}
Define the canonical height by the convention
\begin{equation}\label{eq:height-convention}
 \widehat h_E(P)=\frac12\lim_{n\to\infty}4^{-n}
 h\bigl(x([2^n]P)\bigr),\qquad
 h(a/b)=\log\max\{|a|,b\}.
\end{equation}
Thus $h(x(P))=2\widehat h_E(P)+O_E(1)$, by the Weil--canonical height comparison \cite{SilvermanHeight}.  Let $P_1,\ldots,P_r$ be a basis of the free part of $E(\Q)$, and set
\[
 \mathrm H=(\langle P_i,P_j\rangle)_{i,j},\qquad
 \langle P,Q\rangle=\tfrac12(\widehat h_E(P+Q)-\widehat h_E(P)-\widehat h_E(Q)),
\quad \Reg_E=\det\mathrm H.
\]
For $r>0$, this matrix is positive definite.  Put $t_E=|E(\Q)_{\mathrm{tors}}|$ and let $\vol(B_r)=\pi^{r/2}/\Gamma(1+r/2)$ be the volume of the unit ball in $\R^r$.

\begin{theorem}[Unconditional rational-square height asymptotic]\label{thm:height}
Suppose $r\ge1$.  With
\[
 N_{\square}(X)=\#\{v\in\calV:H(v)\le X,\ D_v=1\},
\]
one has
\begin{equation}\label{eq:squarecount}
 N_{\square}(X)\sim
 \mathrm C_E(\log X)^{r/2},\qquad
 \mathrm C_E=\frac{t_E\,\vol(B_r)}{2^{1+r/2}\sqrt{\Reg_E}}.
\end{equation}
The constant is independent of the finite local-square set $S$.
\end{theorem}
\begin{proof}
By the Mordell--Weil theorem, every point has the form
$T+\sum_i n_iP_i$, with $T$ torsion and $\mathbf n\in\Z^r$.  Its canonical height is $\mathbf n^{\mathrm t}\mathrm H\mathbf n$.  Therefore the number of rational points with canonical height at most $R$ is
\[
 t_E\#\{\mathbf n\in\Z^r:\mathbf n^{\mathrm t}\mathrm H\mathbf n\le R\}
 \sim t_E\frac{\vol(B_r)}{\sqrt{\Reg_E}}R^{r/2},
\]
by elementary lattice-point counting in an expanding ellipsoid.  Since
$\log H(a,b)=h(a/b)+O(1)=2\widehat h_E(P)+O_E(1)$,
the height cutoff corresponds to $R=\frac12\log X+O_E(1)$.  Lemma \ref{lem:points} divides the count by two, except for finitely many torsion and two-torsion exceptions.  This gives \eqref{eq:squarecount}.
\end{proof}

\begin{theorem}[Abelian height singularity]\label{thm:height-zeta}
For $r\ge1$ define
\begin{equation}\label{eq:hsquare}
 \mathcal H_E(s)=\sum_{\substack{v\in\calV\\D_v=1}}H(v)^{-2s}.
\end{equation}
The series converges absolutely and locally uniformly for $\Re s>0$, and along the positive real axis,
\begin{equation}\label{eq:heightpole}
 \boxed{\displaystyle
 \mathcal H_E(s)\sim
 \mathrm C_E\,\Gamma(1+r/2)(2s)^{-r/2}
 \qquad(s\downarrow0).}
\end{equation}
For $r=0$, the square-slope set is finite and $\mathcal H_E(s)$ has a finite limit as $s\downarrow0$.
\end{theorem}
\begin{proof}
Let $A(t)=N_{\square}(e^t)$.  Abel summation gives
\[
 \mathcal H_E(s)=\int_{[0,\infty)}e^{-2st}\,dA(t)
 =2s\int_0^\infty e^{-2st}A(t)\,dt+O(1).
\]
For $r>0$, Theorem \ref{thm:height} implies
$A(t)\sim\mathrm C_Et^{r/2}$.  Substitute $u=2st$ and apply dominated convergence, using the elementary polynomial bound on $A(t)$, to obtain \eqref{eq:heightpole}.  The same bound proves normal convergence for $\Re s>0$.  When $r=0$, all rational points are torsion, so Lemma \ref{lem:points} gives a finite set.
\end{proof}

\begin{remark}[What ``pole'' means]\label{rem:notmero}
Equation \eqref{eq:heightpole} is an \emph{Abelian real-axis pole-type asymptotic}.  It does \emph{not} prove a meromorphic continuation of $\mathcal H_E(s)$ through $s=0$.  Even a Gaussian theta series of a positive-definite lattice has exponentially small Poisson-dual terms of the form $e^{-c/s}$, which preclude a naive meromorphic interpretation.  Nor does the cumulative asymptotic \eqref{eq:squarecount} alone imply a sharp asymptotic in every fixed multiplicative annulus $X\le H<2X$: subtracting two main terms may be smaller than the available lattice-point error.
\end{remark}

\section{A two-variable quadratic-character series}
Define
\begin{equation}\label{eq:two-variable}
 \mathcal Z_E(s,w)=\sum_{v\in\calV}H(v)^{-2s}L(w,\chi_{D_v}).
\end{equation}
For $\Re s>1$ and $\Re w>1$ the series is absolutely convergent, since $|L(w,\chi_D)|\le\zeta(\Re w)$ and the number of primitive slopes of height at most $X$ is $O(X^2)$.  Standard elementary estimates on partial sums of Dirichlet characters and the even quadratic functional equation give normal convergence near $w=1$ for every fixed $\Re s>1$, after the principal-character pole is removed.  They also give normal convergence near $w=0$ when $\Re s>2$.

\begin{proposition}[Two exact boundary specializations]\label{prop:twoslices}
With the indicated domains,
\begin{align}
 \mathcal Z_E(s,0)&=-\tfrac12\mathcal H_E(s)
       &&(\Re s>2),\label{eq:at-zero}\\
 \Res_{w=1}\mathcal Z_E(s,w)&=\mathcal H_E(s)
       &&(\Re s>1).\label{eq:at-one}
\end{align}
The right-hand side of \eqref{eq:at-zero} continues as a separately defined holomorphic function of $s$ for $\Re s>0$.  Neither identity asserts a two-variable continuation to $(s,w)=(0,0)$ or $(1,0)$.
\end{proposition}
\begin{proof}
For a nontrivial primitive even quadratic character, $L(0,\chi_D)=0$.  For $D=1$, $L(0,\chi_1)=\zeta(0)=-1/2$.  Normal convergence near zero proves \eqref{eq:at-zero}.  At $w=1$, every nonprincipal $L(w,\chi_D)$ is holomorphic, while $\zeta(w)$ has residue one.  Local uniform convergence of the nonprincipal contribution gives \eqref{eq:at-one}.  For completeness, the elementary partial-sum estimate $|\sum_{n\le T}\chi_D(n)|\le\min(T,D)$ gives $L(w,\chi_D)\ll_\delta D^\delta(1+\log D)$ uniformly for $|w-1|\le\delta<1/2$.  Since $D_v\ll_EH(v)^4$, any $\Re s>1$ allows a sufficiently small $\delta$ for normal convergence. For $|w|\le\delta$, the even quadratic functional equation and the same bound at $1-w$ give $L(w,\chi_D)\ll_\delta D^{1/2+2\delta}(1+\log D)$. Summing over dyadic height intervals proves normal convergence for $\Re s>2+4\delta$ and gives the claimed neighborhood of zero whenever $\Re s>2$.
\end{proof}

\begin{proposition}[The nonprincipal derivative]\label{prop:class-number}
For $\Re s>2$,
\begin{equation}\label{eq:wderivative}
 \frac{\partial\mathcal Z_E}{\partial w}(s,0)
 =\zeta'(0)\mathcal H_E(s)
 +\frac12\sum_{\substack{v\in\calV\\D_v>1}}
 H(v)^{-2s}\sqrt{D_v}\,L(1,\chi_{D_v}).
\end{equation}
If $h_D$ and $R_D$ denote the class number and regulator of the real quadratic field of fundamental discriminant $D>1$, then each nonsquare summand has positive coefficient
\begin{equation}\label{eq:class-number}
 L'(0,\chi_D)=\tfrac12\sqrt D\,L(1,\chi_D)=h_DR_D>0.
\end{equation}
\end{proposition}
\begin{proof}
The completed even primitive quadratic $L$-function is invariant under $w\mapsto1-w$.  Comparing its leading terms at $0$ and $1$ gives $L'(0,\chi_D)=\sqrt D\,L(1,\chi_D)/2$.  The real quadratic class-number formula gives the final equality \cite{DirichletClass,HeckeTate}.  The conductor bound $D_v\ll_EH(v)^4$ and the standard $L(1,\chi_D)\ll\log(2D)$ ensure absolute convergence for $\Re s>2$.
\end{proof}

\section{The independent bulk limit and its Euler product}
Let
\[
 \Omega_\infty=\{(x,y)\in\R^2:x^2+y^2\le1,\ y>0,\ G(x,y)>0\}.
\]
For $p\in S$, let $\alpha_{p,S}$ be the conditional Haar probability that $G(a,b)$ is a nonzero $p$-adic square, given that $(a,b)\in\Z_p^2$ is primitive.  Put
\begin{equation}\label{eq:density}
 \mathrm d_{E,S}=\frac{\operatorname{area}(\Omega_\infty)}{\zeta(2)}
 \prod_{p\in S}\alpha_{p,S}>0.
\end{equation}
For positivity at a prime $p$, take $a=1$ and $b=Ap^{2n}$ for large $n$. Then $G(1,b)=A^2p^{2n}(1+O_E(p^{2n}))$ in $\Q_p$, where the parenthesis tends $p$-adically to one and is a square for large $n$, also at $p=2$. Continuity gives an open set of primitive local solutions. Over $\R$ take $a$ of the same sign as $A$ and $b>0$ small, then rescale into the unit disk. Elementary lattice counting with Moebius inversion and the Chinese remainder theorem, approximating the boundaries of $p$-adic square loci by finite congruences, gives
\begin{equation}\label{eq:allcount}
 N(X):=\#\{v\in\calV:H(v)\le X\}
 \sim\mathrm d_{E,S}X^2.
\end{equation}

For each odd good prime $p\notin S$, let
\[
 a_p=p+1-\#E(\F_p),\qquad
 b_p=\#\{x\in\mathbb P^1(\F_p):G(x)=0\}.
\]
The latter includes the branch point at infinity and equals the number of Frobenius-fixed geometric branch points.  In particular $1\le b_p\le4$.  Define
\begin{equation}\label{eq:localmoments}
 \mu_p=-\frac{a_p}{p+1},\qquad
 \delta_p=\frac{p b_p}{(p+1)^2}.
\end{equation}

\begin{lemma}[Exact good-prime character average]\label{lem:local}
On the space of primitive $p$-adic slopes with normalized Haar measure,
\[
 \mathbb E[\chi_{D_v}(p)]=\mu_p,\qquad
 \mathbb P(p\mid D_v)=\delta_p,
\]
and, with $t=p^{-w}$,
\begin{equation}\label{eq:Mp}
 \mathbb E[(1-\chi_{D_v}(p)t)^{-1}]
 =\frac{1+\mu_pt-\delta_pt^2}{1-t^2}.
\end{equation}
At $p\in S$, the conditional local factor is $(1-p^{-w})^{-1}$.
\end{lemma}
\begin{proof}
The $p+1$ residue classes of primitive projective slopes are equiprobable.  The point-count identity for the double cover $y^2=G$ is
$\sum_{x\in\mathbb P^1(\F_p)}(G(x)/p)=-a_p$, giving the first expectation.  A simple branch root has probability $1/(p+1)$ at the first residue level.  Conditional on that root, the probability that the valuation of $G$ is odd is
\[
 \sum_{j\ge0}(p-1)p^{-(2j+1)}=\frac{p}{p+1}.
\]
This gives $\delta_p$.  The even-valuation branches have zero mean quadratic character after averaging their unit residues.  Hence all positive odd moments equal $\mu_p$, while all positive even moments equal $1-\delta_p$.  Summing the geometric series proves \eqref{eq:Mp}.  At primes in $S$, the local squareclass is trivial, so $\chi_{D_v}(p)=1$.
\end{proof}

\begin{theorem}[Bulk mean and Abelian boundary residue]\label{thm:bulk}
For $\Re w>1$, the limit
\begin{equation}\label{eq:bulkmean}
 \mathcal M_{E,S}(w)=\lim_{X\to\infty}
 \frac{1}{N(X)}\sum_{\substack{v\in\calV\\H(v)\le X}}
 L(w,\chi_{D_v})
\end{equation}
exists and is the absolutely convergent Euler product
\begin{equation}\label{eq:eulerproduct}
 \boxed{\displaystyle
 \mathcal M_{E,S}(w)=
 \prod_{p\in S}\frac1{1-p^{-w}}
 \prod_{p\notin S}
 \frac{1-\frac{a_p}{p+1}p^{-w}
       -\frac{p b_p}{(p+1)^2}p^{-2w}}
      {1-p^{-2w}}.}
\end{equation}
For real $w>1$,
\begin{equation}\label{eq:abelian-residue}
 \lim_{s\downarrow1}(s-1)\mathcal Z_E(s,w)
 =\mathrm d_{E,S}\mathcal M_{E,S}(w).
\end{equation}
The limit in \eqref{eq:abelian-residue} is an Abelian boundary residue; a meromorphic continuation of $\mathcal Z_E(s,w)$ across $s=1$ is not asserted.  If such a continuation with a simple pole exists, its residue must equal the displayed expression.
\end{theorem}
\begin{proof}
For any fixed integer $n$, $\chi_{D_v}(n)$ is a product of bounded local functions at the primes dividing $n$.  They are locally constant outside zero sets of $G$, which have $p$-adic measure zero.  Equidistribution of primitive lattice points under the finite local restrictions and the Chinese remainder theorem show that the mean of $\chi_{D_v}(n)$ tends to the product of its local means.  Since for $\Re w>1$,
$L(w,\chi_D)=\sum_{n\ge1}\chi_D(n)n^{-w}$ is absolutely bounded by $\zeta(\Re w)$, dominated convergence permits termwise passage to the limit.  Local factorization and Lemma \ref{lem:local} give \eqref{eq:eulerproduct}.  Finally the weighted counting function is asymptotic to
$\mathrm d_{E,S}\mathcal M_{E,S}(w)X^2$.  Partial summation of $H^{-2s}$ proves \eqref{eq:abelian-residue}.
\end{proof}

\subsection{Meromorphic continuation of the bulk coefficient}
Let
\begin{equation}\label{eq:branchalgebra}
 \calB=\Q\times\Q[T]/(f(T)),\qquad
 b=\#\{\text{irreducible factors of }G\text{ in }\Q[T,U]\}.
\end{equation}
Since $f$ is separable, $\calB$ is an etale $\Q$-algebra of degree four, and $\zeta_{\calB}(z)$ is the product of its Dedekind zeta functions.  Its pole at $z=1$ has order $b$; equivalently, $b$ is the number of rational Galois orbits among the four branch points.  If $f$ is irreducible, $b=2$.

A superscript $S$ denotes deletion of Euler factors at $S$.  The omitted elliptic Euler factors are nonzero at $s=1$, so $L^S(E,s)$ and the full $L(E,s)$ have the same central vanishing order.

\begin{theorem}[The character bulk has order $1+b-a_E$]\label{thm:order}
The Euler product $\mathcal M_{E,S}(w)$ of Theorem \ref{thm:bulk} has a meromorphic continuation to a neighborhood of $w=0$ of the form
\begin{equation}\label{eq:factorization}
 \boxed{\displaystyle
 \mathcal M_{E,S}(w)=
 U_{E,S}(w)\,
 \frac{\zeta(2w)}
 {L^S(E,1+w)\,\zeta^S(1+2w)\,
  \zeta_{\calB}^S(1+2w)},}
\end{equation}
where $U_{E,S}(w)$ is holomorphic and nonzero near $w=0$.
Consequently
\begin{equation}\label{eq:order}
 \boxed{\displaystyle
 \ord_{w=0}\mathcal M_{E,S}(w)=1+b-a_E.}
\end{equation}
This statement uses modularity of $E$, not BSD or finiteness of the Tate--Shafarevich group.
\end{theorem}
\begin{proof}
For $p\notin S$, set $t=p^{-w}$ and
\[
 V_p(w)=1-\frac{a_p}{p+1}t-\frac{p b_p}{(p+1)^2}t^2,
 \qquad L_p(E,1+w)^{-1}=1-\frac{a_p}{p}t+\frac1p t^2.
\]
The permutation representation on the four branch points has Frobenius trace $b_p$, so
\[
 \zeta_{\calB,p}(1+2w)^{-1}
 =1-\frac{b_p}{p}t^2+O(p^{-2}|t|^4).
\]
Multiplying this with $1-p^{-1}t^2$ and $L_p(E,1+w)^{-1}$ gives
\[
 V_p(w)=L_p(E,1+w)^{-1}(1-p^{-1}t^2)
 \zeta_{\calB,p}(1+2w)^{-1}
 \bigl(1+O(p^{-3/2+4\epsilon})\bigr)
\]
uniformly for $\Re w\ge-\epsilon$, with $0<\epsilon<1/8$. Indeed $|t|\le p^\epsilon$ throughout this half-plane, and the displayed local comparison factors have no zeros there: the reciprocal roots for the elliptic factor have absolute value $\sqrt p$ in the unshifted Euler variable, and those for the permutation representation have absolute value one. Hasse's bound $|a_p|\le2\sqrt p$ and $b_p\le4$ give the stated summable error. The product of the local correction factors therefore converges normally on $\Re w>-\epsilon$. This common domain reaches the original half-plane $\Re w>1$, so the factorization gives analytic continuation from that Euler product.

We must also check that it is nonzero there.  For $p\ge11$, Hasse's lower bound
$\#E(\F_p)\ge p+1-2\sqrt p>5\ge1+b_p$
ensures a nonbranch point with nonzero square value of $G$.  Hence
$V_p(0)=1+\mu_p-\delta_p>0$.  All comparison factors at $w=0$ are positive and nonzero.

Finally,
\[
 \prod_{p\in S}(1-p^{-w})^{-1}\zeta^S(2w)
 =\zeta(2w)\prod_{p\in S}(1+p^{-w}).
\]
The finite product is a holomorphic unit at zero. This proves
\eqref{eq:factorization}. The factors $\zeta^S(1+2w)$ and
$\zeta_{\calB}^S(1+2w)$ have poles of orders one and $b$ respectively.
Since $\zeta(0)=-1/2$ and $L^S(E,1+w)$ has order $a_E$,
\eqref{eq:order} follows.
\end{proof}

\section{Deleting the square slopes and differentiating the mean}
\begin{proposition}\label{prop:delete}
For fixed $\Re w>1$, deleting the slopes for which $D_v$ belongs to any
fixed finite set of fundamental discriminants leaves the limit
$\mathcal M_{E,S}(w)$ unchanged.
\end{proposition}
\begin{proof}
Theorem \ref{thm:height} bounds the number of square slopes of height
at most $X$ by $O((1+\log X)^{r/2})$. For fixed $D>1$, the slopes with
$D_v=D$ lift to rational points on the quadratic twist $f(x)=Dy^2$.
Applying the same height comparison gives the bound
$O_D((1+\log X)^{r_D/2})$, with $r_D$ the Mordell--Weil rank of the twist.
Each slope in that set contributes the same value $L(w,\chi_D)$.
A finite sum of these bounds is $o(N(X))$ by \eqref{eq:allcount}.
\end{proof}

Write
\[
 M_X(w)=\frac1{N(X)}\sum_{\substack{v\in\calV\\ H(v)\le X}}L(w,\chi_{D_v}).
\]
At zero, $M_X(0)=-N_\square(X)/(2N(X))$. For $r>0$ it is asymptotic to
$-\mathrm C_E(\log X)^{r/2}/(2\mathrm d_{E,S}X^2)$.
\begin{proposition}
For every curve in the given setting, $M_X'(0)\to+\infty$ as
$X\to\infty$. In particular the finite-height means cannot converge
locally uniformly near zero to a holomorphic continuation of the bulk mean.
\end{proposition}
\begin{proof}
By \eqref{eq:class-number},
\[
 M_X'(0)=\frac{N_\square(X)}{N(X)}\zeta'(0)+
 \frac1{N(X)}\sum_{\substack{v\in\calV,\ H(v)\le X\\D_v>1}}h_{D_v}R_{D_v}.
\]
Proposition \ref{prop:delete} shows that $D_v$ tends to infinity in
probability under the uniform measure on the slopes. A fundamental real
quadratic unit $\epsilon>1$ satisfies $\epsilon\ge\sqrt D/2$ (subtract
its conjugate and use $|\epsilon'|<1$). Hence
$h_DR_D\ge\tfrac12\log D-\log2$. The nonprincipal terms are nonnegative,
and their mean tends to infinity. The principal term tends to zero.
Local uniform convergence of holomorphic functions would give
convergence of their first derivatives.
\end{proof}

Even for a meromorphic function of two variables, different orders of taking
a residue and a specialization can give different results. For example
\[
 F(x,w)=\frac1{x(x+w)}-\frac1{x^2}
\]
satisfies $F(x,0)=0$, while $\Res_{x=0}F(x,w)=1/w$ for $w\ne0$.
For the actual mean, Proposition \ref{prop:delete} goes further.
Removing every global rational-square slope preserves the bulk coefficient
in its initial half-plane and hence preserves its meromorphic continuation.

\section{Retaining the square part}
Write $G(v)=D_v q_v^2$ with $q_v>0$ rational, and put
$A(v)=G(v)/H(v)^4$. For a bounded smooth angular weight $\varphi$ define
\[
 \mathcal Z_E(s,w,\tau;\varphi)=
 \sum_{v\in\calV} H(v)^{-2s}q_v^{-\tau}
 \varphi(v/H(v))L(w,\chi_{D_v}).
\]
The even quadratic functional equation has gamma factor
$\gamma(w)=\pi^{w-1/2}\Gamma((1-w)/2)/\Gamma(w/2)$. Substitution of
$D_v=H(v)^4 A(v)/q_v^2$ in each summand yields
\[
 \mathcal Z_E(s,w,\tau;\varphi)=\gamma(w)
 \mathcal Z_E(s-1+2w,1-w,\tau+1-2w;\varphi A^{1/2-w}).
\]
The weight on the right is evaluated on the positive arcs $A>0$ and may grow at their endpoints. On the sampled slopes $H^{-4}\le A\ll_E1$, so for fixed $w$ its growth is at most a fixed power of $H$. The identity holds in common domains of absolute convergence, with both sides
interpreted meromorphically at the exceptional character parameters when
justified. For fixed $w,\tau$ away from poles such domains exist for
sufficiently large $\Re s$, since $G$ is an integer of size $O_E(H^4)$,
$D_v\ll_E H^4$, and $q_v$ has denominator at most two. Both $s+\tau$
and $s+w-1$ are invariant under the reflection. On square slopes,
$q_v=H(v)^2\sqrt{A(v)}$ exactly, so the square-part variable retains
data lost in the unweighted bulk. The reflection still supplies no
comparison of the two invariant directions.

\section{Two quartics}\label{sec:examples}
For $E_4:y^2+xy=x^3-x^2-79x+289$, put $u=4x-1$, $v=8y+4x$.
Then $v^2=u^3-1267u+17230$, giving
\[
 G_4(a,b)=b(a^3-1267ab^2+17230b^3).
\]
Direct substitution gives $G_4(15,1)=40^2$ and
$G_4(5,3)=4785\cdot16^2$. The cubic is irreducible modulo three,
so the branch algebra has $b=2$ rational components.
For $E_2:y^2+y=x^3+x^2-2x$, completing the square instead gives
\[
 G_2(a,b)=b(4a^3+4a^2b-8ab^2+b^3).
\]
Here $G_2(3,1)=11^2$ and $G_2(-8,5)=4785$. The rational-root test
shows the cubic has no rational root, hence $b=2$ again. These
examples of $D=4785$ concern the unrestricted quartics. The local
square restrictions at $3$ and $5$ exclude that discriminant in $\calV$.

Taking the independently recorded ranks two and four for $389.a1$ and
$234446.a1$ gives, respectively, growth proportional to $s^{-1}$ and
$s^{-2}$ in Theorem \ref{thm:height-zeta} \cite{curves}. The bulk formula
gives $3-a_E$ on both curves. The height formula concerns the actual
rational points. The analytic order concerns a coefficient preserved
after those points are deleted. Identifying the two ranks would require
an additional arithmetic argument.

\section{A self-dual prime sum and deletion of rational rays}\label{sec:ray-deletion}

For a good odd prime $p\notin S$ put $\varphi_p(v)=(G(v)/p)$, using the
Legendre symbol extended by zero. Homogeneity makes this constant on the
nonzero vectors on each line in $\F_p^2$. The double cover of the
projective line defined by the quartic has $p+1+\sum_{[v]}\varphi_p(v)$
rational points, so $\sum_{[v]}\varphi_p(v)=-a_p$. For
$J(u_1,u_2)=(-u_2,u_1)$, put $e_p(x)=\exp(2\pi ix/p)$. Summing first
along each nonzero line gives the finite Fourier identity
\[
 \sum_{v\bmod p}\varphi_p(v)e_p(u\cdot v)
 =p\varphi_p(Ju)+a_p-p a_p\,\mathbf 1_{u=0}.
\]
Indeed the character sum on a line with its zero removed is either
$p-1$ or $-1$, according as the line is perpendicular to $u$ or not.

Take $W(x)=(1-2\pi\|x\|^2)e^{-\pi\|x\|^2}$ on $\R^2$ and the Fourier
transform with kernel $e^{-2\pi i x\cdot y}$. Then $\widehat W=-W$.
Poisson summation on each residue class of $\Z^2$ modulo $p$, followed
by the finite identity and the ordinary Gaussian Poisson identity
$\sum_m W(m/\sqrt p)=-p\sum_m W(\sqrt p\,m)$, gives exactly
\begin{equation}\label{eq:ray-prime}
 S_p:=\sum_{v\in\Z^2}\varphi_p(v)W(v/\sqrt p)
       =a_p\sum_{m\in\Z^2}W(\sqrt p\,m).
\end{equation}
The rightmost sum equals $1+O(p e^{-\pi p})$. The prime series
$\sum_p S_p\log p\, p^{-s}$ consequently has the same continued
residue at $s=1$ as the prime-only series for $a_p$, namely
$\tfrac12-a_E$. To check the half, remove prime powers from
$-(L^S)'/L^S$. The second-power coefficient is $a_p^2-2p$ and its
series has residue $-1/2$ there, by the Rankin--Selberg simple pole
and regularity and nonvanishing of the adjoint factor at the edge (also for CM). Equivalently the squared normalized Hecke traces have prime average one; substituting $2s-1$ into the logarithmic derivative supplies the factor $1/2$.
Powers three and higher converge normally near $s=1$.

For a primitive, unoriented rational square ray $v$, let $h=H(v)$ and
$g(t)=(1-2\pi t^2)e^{-\pi t^2}$. Its full contribution to
\eqref{eq:ray-prime} is zero at primes dividing $G(v)$, and otherwise
equals
\begin{equation}\label{eq:ray-line}
 K(h/\sqrt p)-K(h\sqrt p),\qquad K(t)=\sum_{n\in\Z}g(nt).
\end{equation}
Here $p\mid n$ gives zero in the prime sum. The integral of $g$ on
$\R$ is zero. One-dimensional Poisson summation gives $K(t)=O_A(t^A)$
for $0<t\le1$ and $K(t)=1+O_A(t^{-A})$ for $t\ge1$, for every $A>0$.
For each fixed ray the contribution is thus $-1+O_{v,A}(p^{-A})$.
Deleting that one ray adds one to the continued prime residue. For
any fixed radial Schwartz function with Fourier transform $-W$ the
same proof gives $-W(0)$ in place of $-1$: Fourier slicing identifies
its line integral with the integral of $\widehat W$ on a perpendicular
line, which equals its negative by radiality. The coefficient of
$a_p$ in \eqref{eq:ray-prime} also tends to $W(0)$.

For the full set of rational square rays and $r\ge1$ one obtains
\begin{equation}\label{eq:ray-slice}
 S_{p,\square}
 =-\mathrm C_E\left(\tfrac12\log p\right)^{r/2}
      +O_E\!\left(1+(\log p)^{(r-1)/2}\right).
\end{equation}
Here is a uniform justification. The ellipsoid count used in
Theorem~\ref{thm:height} has error $O(1+T^{(r-1)/2})$ in a
logarithmic radius cutoff $T=\log H$. Cover lattice points by their
unit fundamental cells; only cells in a boundary layer of bounded
Euclidean thickness in the Mordell--Weil lattice can disagree with
volume. The bounded error between $2\widehat h$ and $\log H$ is
absorbed by the same bound. In every shell $j\le\log h<j+1$ this
gives $O_E(1+j^{(r-1)/2})$ rays. Put $T=\tfrac12\log p$. Equation
\eqref{eq:ray-line} differs from $-\mathbf1_{\log h\le T}$ by
$O_A(e^{-A|\log h-T|}+ (h\sqrt p)^{-A})$. Summing this shell
error costs $O_E(1+T^{(r-1)/2})$.

It remains to exclude $p\mid G(v)$ in the main count $h\le\sqrt p$.
The point underlying such a ray reduces to two-torsion or to the
origin, so $2P$ reduces to the origin. If $2P\ne O$, the denominator
of $x(2P)$ is divisible by $p^2$ on a fixed integral short
Weierstrass model away from finitely many primes. Thus
$8\widehat h(P)+O_E(1)=h(x(2P))\ge2\log p$.
Apply this also to differences of points in the reduction kernel
of $[2]$. In the free Mordell--Weil lattice, distinct such points
modulo torsion are separated by $\gg_E\sqrt{\log p}$ in the height
norm. The points under consideration lie in a ball of radius
$O_E(\sqrt{\log p})$, so packing bounds their number by $O_E(1)$.
Finitely many torsion points have bounded contribution. This proves
\eqref{eq:ray-slice}.

In particular for $234446.\mathrm{a}1$ deleting square rays outside
any fixed rank-two subgroup increases $S_p$ by
$\mathrm C_E(\log p)^2/4+O_E((\log p)^{3/2})$. The subgroup itself
contributes $O_E(\log p)$ by the same count. This distinguishes
subgroups at every sufficiently large prime. It does not determine
the central residue: deletion of even one more fixed ray changes that
residue by one and leaves all leading powers in \eqref{eq:ray-slice}
unchanged.
\section{Local probabilities on two integral quartic charts}

At $2$ the measure on slopes depends on coordinates. For this section
we use the original abscissa $x$ of each minimal Weierstrass equation
and complete the square in $y$, giving
\[
 F_4(x)=4x^3-3x^2-316x+1156,\qquad F_2(x)=4x^3+4x^2-8x+1
\]
for $234446.\mathrm{a}1$ and $389.\mathrm{a}1$ respectively.
Write $G_i^0(m,n)=n^4 F_i(m/n)$. In Section~\ref{sec:examples},
$G_4$ instead uses the abscissa $4x-1$. That change of variable
does not preserve uniform projective measure at $2$.

For $v=2^eu$ with $u$ odd, $\chi_D(2)$ equals $+1$ or $-1$
when $e$ is even and $u$ is respectively $1$ or $5\bmod8$; it is
zero otherwise. Uniform projective slopes have two charts. Those
with $n$ odd have mass $2/3$, with $m/n$ uniform in $\Z_2$.
The other chart has mass $1/3$, with $n/m$ uniform in $2\Z_2$.
The conditional distributions in the order $(+1,-1,0)$ are
\[
\begin{array}{c|c|c|c}
 & n\text{ odd}& n\text{ even}&\text{overall}\\ \hline
E_4&(13/16,1/16,1/8)&(1/24,1/24,11/12)&(5/9,1/18,7/18)\\
E_2&(1,0,0)&(1/12,1/12,5/6)&(25/36,1/36,5/18)
\end{array}
\]

Indeed $F_2(x)\equiv1\bmod8$ for every integral $x$. For $F_4$,
odd $x$ gives $F_4(x)\equiv1\bmod8$, and $4\mid x$ gives
$F_4(x)/4\equiv1\bmod8$. If $x=2y$ with $y$ odd, then
$F_4(2y)/4=8y^3-3y^2-158y+289$ takes residues $8,4,8,20\bmod32$
as $y$ runs through $1,3,5,7\bmod8$. These facts give the
odd-denominator column.

On the other chart write $t=n/m=2^ku$, with $u$ uniform and odd.
Here $k\ge1$ has probability $2^{-k}$. The reversed quartics are
\[
 G_2^0(1,t)=t(4+4t-8t^2+t^3),\qquad
 G_4^0(1,t)=t(4-3t-316t^2+1156t^3).
\]
For $E_2$ the parenthesis has valuation two. Its quotient by four
is $5\bmod8$ when $k=1,2$, and $1\bmod8$ when $k\ge3$.
Only even $k$ contribute a sign, each sign with conditional
probability $1/4$. For $E_4$ at $k=1$ the normalized odd part
is $3$ or $7\bmod8$. At $k=2$ the value equals $16uH(u)$ where
$H(u)=1-3u-1264u^2+18496u^3$. The map $u\mapsto uH(u)$ maps
odd residues bijectively to even residues modulo every $2^j$ by
Hensel lifting: its derivative is odd throughout the chart.
Each nonzero sign then has conditional probability
$\tfrac14\sum_{j\ge1}2^{-2j}=1/12$. At $k\ge3$ the parenthesis has
valuation two; at $k=4$ the normalized odd part is $u+4\bmod8$,
and for even $k\ge6$ it is $u\bmod8$. Summing with weights
$2^{-k}$ gives the even-denominator column.

At the odd bad primes the change $x\mapsto4x-1$ preserves the
measure and the quadratic squareclass. Direct reduction of $F_i$
gives $4(x-\alpha)^2(x-\beta)$, with
\[
 \begin{array}{c|r|r|r|r}
 &p&\alpha&\beta&F_i(\alpha)/p\bmod p\\ \hline
 E_4&117223&7592&72734&43107\\
 E_2&389&299&179&369
 \end{array}
\]
Thus the double-root tube has valuation exactly one and contributes
zero. The two simple-root directions are $\beta$ and infinity.
Each contributes $1/[2(p+1)^2]$ to the projective squareclass mass
by summing even positive valuations in a simple Hensel tube.
In both rows $\alpha-\beta$ is a nonzero square modulo $p$.
Away from these three directions the quadratic-character sum is
$-1$, so exactly $(p-3)/2$ directions have square value. Hence
\[
 \rho_p=\frac{p-3}{2(p+1)}+\frac1{(p+1)^2},\qquad
 \rho_{117223}=\frac{6870498641}{13741466176},\quad
 \rho_{389}=\frac{75271}{152100}.
\]
All these numbers describe the indicated unrestricted projective
charts before imposing the local square restrictions of $\calV$.

There is a related transform for the square indicator (including
zero values of $G$). At an odd prime let $s_p(v)$ indicate this
condition for $v\ne0$ and put $s_p(0)=0$. Write $t_p$ for the number
of allowed projective residue directions. With Fourier normalization
$p^{-2}$ the zero coefficient is $(p-1)t_p/p^2$, and each nonzero
frequency has coefficient $(p\,s_p(Ju)-t_p)/p^2$. Summing first
over the $p-1$ nonzero scalar representatives of a projective line
proves the formula; the Chinese remainder theorem gives all
squarefree-modulus versions.

\subsection{An upper bound on a shell derivative}
Use $G^0_4$ and the max-norm cutoff
$H\le\max(|u|,v)<2H$, $v>0$, $\gcd(u,v)=1$. Retain positive
$G^0_4$ with fundamental discriminant $D>1$, optionally imposing
further local square restrictions. For an integer $H\ge1$ put
\[
 M_H(w)=H^{-4+4w}\sum_{\rm retained}L(w,\chi_D).
\]
Each character is even and nonprincipal, so $M_H(0)=0$ and the
Dirichlet functional equation gives
\[
 0\le M'_H(0)=\frac1{2H^4}\sum_{\rm retained}\sqrt D\,L(1,\chi_D)
 \le4\sqrt{94656}\bigl(\log(94656H^4)+2\bigr).
\]
Indeed there are fewer than $8H^2$ pairs and
$D\le4G_4^0\le94656H^4$. A sum through $D$ and partial summation
using the character period bound the Dirichlet value by $\log D+2$.
The dyadic sum therefore satisfies
\[
 \limsup_{s\downarrow0}s^2\sum_{j\ge0}2^{-js}M'_{2^j}(0)
 \le\frac{16\sqrt{94656}}{\log2}.
\]
This bound is a finite-height estimate. For example for any constant
$K$, $K(1-H^{-w})^3/w+wH^{-w}$ has derivative $1$ at zero for each
finite $H$, but its limit at fixed real $w>0$ is $K/w$.

\section{Signed degree-two divisors and rational heights}
Let $\sigma:X=\operatorname{Sym}^2 E\to E$ be the sum map. Translating
both members of an unordered pair by $-T$ when its sum is $2T$ gives
an isomorphism over $\Q$
\[
 E\times\mathbb P^1\xrightarrow{\ \sim\ } X\times_{E,[2]} E,\qquad
 (T,x(P))\longmapsto(\{T+P,T-P\},T).
\]
The inverse uses $E/\{\pm1\}\simeq\mathbb P^1$. By the projective
bundle formula, every line bundle on this product is obtained from a
line bundle on $E$ and $\mathcal O_{\mathbb P^1}(d)$. Its Weil height
therefore separates up to a bounded term. The identity for canonical
heights is exact:
\[
 \widehat h(T+P)+\widehat h(T-P)=2\widehat h(T)+2\widehat h(P).
\]

\subsection{Frobenius traces}
Let $\pi:E^2\to X$ be the quotient by the switch $\tau$ and put
$\mathcal F=(\pi_*\Q_\ell)^-$ at a good prime $p\ne\ell$.
On pairs the trace has values $1,-1,0$ for distinct split points,
a quadratic conjugate pair, and a repeated point respectively. If
$q=p^n$, $N_q=q+1-a_q=\#E(\F_q)$, then $N_{q^2}=(q+1)^2-a_q^2$.
Hence the total signed trace and its mean on $X(\F_q)$ are
\[
 \frac{N_q^2-N_{q^2}}2=-a_q N_q,\qquad -\frac{a_q}{q+1}.
\]
The denominator in the mean comes from the $q+1$ points of each
projective-line fibre. For $V=H^1(E_{\overline{\F}_p},\Q_\ell)$ one has
$R\sigma_*\mathcal F\simeq V[-1]$ as a geometrically constant
complex with its natural Frobenius action. In fact
$(A,B)\mapsto(A+B,A)$ identifies the ordered-pair family with a
product; the involution on the second factor is $A\mapsto S-A$
above $S$. It acts by $-1$ in degree one and by $+1$ in degrees
zero and two. Proper base change and the idempotent $(1-\tau)/2$
therefore give the assertion, including over a fibre with coincident
points. If $\alpha,\beta$ are the two Frobenius eigenvalues on $V$,
the resulting local zeta factor is
\[
 \exp\left(\sum_{n\ge1}\frac{-a_{p^n}N_{p^n}}n u^n\right)=
 \frac{(1-\alpha u)(1-\beta u)(1-p\alpha u)(1-p\beta u)}
 {(1-\alpha^2u)(1-pu)^2(1-\beta^2u)} .
\]
Here $a_{p^n}$ denotes $\alpha^n+\beta^n$ in this subsection.
In contrast, the ordinary generating function of the normalized
means is not rational: expanding $(p^n+1)^{-1}$ gives simple poles
at $p^j/\alpha,p^j/\beta$ for $j\ge1$. Distinct $j$ have different
pole moduli, and coincident poles for $\alpha=\beta$ add. The resulting
series extends meromorphically off these points by normal convergence
on compact subsets. Thus the normalized means cannot themselves be
the power traces of a single finite-dimensional operator.

\subsection{A signed height series}
On $\mathbb P^1(\Q)$ write $H_0$ for primitive max height and put
$\epsilon_E(u)=\#x^{-1}(u)(\Q)-1$. Its three values use the same
splitting convention over $\Q$. For $\Re w>2$,
\[
 \Sigma_E(w):=\sum_{u\in\mathbb P^1(\Q)}\epsilon_E(u)H_0(u)^{-w}
 =\sum_{P\in E(\Q)} H_0(x(P))^{-w}
       -4\frac{\zeta(w-1)}{\zeta(w)} .
\]
Here $x(O)=\infty$. The factor four follows from the count $4\phi(n)$
of primitive projective directions of height $n$, with $\phi(1)=1$.
The right side continues meromorphically to $\Re w>0$.
Write $t_E=|E(\Q)_{\rm tors}|$, $t_2=|E(\Q)[2]|$ and
$\Theta_E(z)=\sum_{T\in E(\Q)} e^{-z\widehat h(T)}$.
On the sector $\sigma(D)\in 2E(\Q)$, count an effective divisor $D$
with the mean of its weights over its $t_2$ rational lifts to the
product above. A lift $(T,u)$ has weight
$\epsilon_E(u)e^{-z\widehat h(T)}H_0(u)^{-w}$. The resulting series
is exactly $\Theta_E(z)\Sigma_E(w)/t_2$.
When $r>0$, lattice counting or Poisson summation, together with
$\log H_0(x(P))=2\widehat h(P)+O_E(1)$, gives
\[
 \frac{\Theta_E(z)\Sigma_E(z)}{t_2}\sim
 \frac{t_E^2\pi^r}{2^{r/2}t_2\Reg_E}z^{-r}\quad(z\downarrow0).
\]
The zeta quotient is regular at zero. For $r=0$ its subtraction gives
$\Sigma_E(0)=t_E-2/3$. For the corresponding positive count of
\emph{distinct split} divisors with $\widehat h(A)+\widehat h(B)\le T$
and $A+B\in2E(\Q)$, the leading term for $r>0$ is
\[
 \frac{t_E^2\pi^r}{2^{r+1}t_2\,r!\Reg_E}T^r .
\]
Indeed the finite-index condition has index $2^r t_2$ in the
ordered-pair group, then switching contributes a factor $1/2$;
the diagonal contributes $O(T^{r/2})$.

\subsection{An independence weight for tuples}
On the free Mordell--Weil lattice $\Gamma$ form the Gaussian sum
\[
 W_k(z)=\sum_{P_1,\ldots,P_k\in\Gamma}
 \det(\langle P_i,P_j\rangle)_{i,j\le k}
 e^{-z\sum_i\widehat h(P_i)},\qquad z>0 .
\]
The Gram determinant vanishes exactly for dependent tuples.
Thus $W_k(z)>0$ exactly for $k\le r$. If $1\le k\le r$, lattice
Riemann sums yield
\[
 W_k(z)\sim k!\binom rk
  \left(\frac{\pi^{r/2}}{2\sqrt{\Reg_E}}\right)^k
  z^{-k(r/2+1)}.
\]
In the Gaussian integral, Cauchy--Binet gives one squared minor for
each set of $k$ coordinates. Its expectation is $k!/2^k$ because
the coordinates under normalized density proportional to
$e^{-\|x\|^2}$ have variance $1/2$. This gives the coefficient with
the height convention of \eqref{eq:height-convention}. The exponent
detects the number of independent arithmetic directions available
to the tuple.

\subsection{Positive Euler coefficients and their rightmost pole}
For a good prime, put $u=\sqrt p$, $X=p^{-s}$. The local factor of
$\zeta(s-\tfrac12)^2/L(E,s)$ has expansion
\[
 \frac{1-a_pX+pX^2}{(1-uX)^2}
 =1+\sum_{n\ge1}n u^{n-1}(2u-a_p)X^n .
\]
All coefficients are nonnegative by Hasse. At the bad primes the
numerator becomes $1-a_pX$, with $a_p\in\{-1,0,1\}$, and the same
nonnegativity holds. The zeta functions introduce a double pole at
$s=3/2$ by the nonvanishing of the normalized cuspidal $L$-function
on $\Re s=1$ \cite{JS}. At the central point the pole order is $a_E$
when $a_E>0$. A positive-coefficient Tauberian estimate from the
rightmost singularity hence describes a different asymptotic.

\subsection{Vertical factors in a model of the symmetric square}
The generic fibre $\operatorname{Sym}^2E$ is a surface, so a regular
arithmetic model has dimension three. Blowing up a regular closed
centre $Z$ of codimension $c$ in such a model $Y$ gives
\[
 \zeta_{\operatorname{Bl}_Z Y}(s)/\zeta_Y(s)
 =\prod_{j=1}^{c-1}\zeta_Z(s-j).
\]
Indeed the exceptional divisor is the projectivized normal bundle,
and the complements of the centre and exceptional divisor agree.
A closed point with residue cardinality $q$ introduces factors
$(1-q^{1-s})^{-1}(1-q^{2-s})^{-1}$. They record the choice of model;
the signed trace and the rank exponent above were obtained on the
generic family and its good reductions.

\raggedright
\begin{thebibliography}{99}
\bibitem{JS}
H.~Jacquet and J.~A.~Shalika, \emph{On Euler products and the
classification of automorphic representations. II},
Amer. J. Math. \textbf{103} (1981), 777--815.
\bibitem{SilvermanHeight}
J. H. Silverman, \emph{The difference between the Weil height and the canonical height on elliptic curves},
Math. Comp. \textbf{55} (1990), 723--743.
\url{https://doi.org/10.1090/S0025-5718-1990-1035944-5}.
\bibitem{SilvermanAEC}
J. H. Silverman, \emph{The Arithmetic of Elliptic Curves}, 2nd ed.,
Graduate Texts in Mathematics 106, Springer, 2009, Chapter VIII.
\bibitem{BCDT}
C. Breuil, B. Conrad, F. Diamond and R. Taylor,
\emph{On the modularity of elliptic curves over $\Q$: wild $3$-adic exercises},
J. Amer. Math. Soc. \textbf{14} (2001), 843--939.
\url{https://doi.org/10.1090/S0894-0347-01-00370-8}.
\bibitem{Deligne}
P. Deligne, \emph{La conjecture de Weil II},
Publ. Math. IH\'ES \textbf{52} (1980), 137--252.
\url{https://doi.org/10.1007/BF02684780}.
\bibitem{DirichletClass}
T. Mitsuhiro, T. Nakahara and T. Uehara, \emph{The class number formula of a real quadratic field and an estimate of the value of a unit},
Canad. Math. Bull. \textbf{38} (1995), 98--103.
\url{https://doi.org/10.4153/CMB-1995-013-6}.
\bibitem{HeckeTate}
J. Tate, \emph{Fourier analysis in number fields and Hecke's zeta-functions},
in \emph{Algebraic Number Theory}, Academic Press, 1967.
\bibitem{Cremona}
J. E. Cremona, \emph{Algorithms for Modular Elliptic Curves}, 2nd ed.,
Cambridge University Press, 1997. \url{https://johncremona.github.io/book/amec.html}.
\bibitem{curves}
The LMFDB Collaboration, elliptic curve records
\href{https://www.lmfdb.org/EllipticCurve/Q/389/a/1}{389.a1} and
\href{https://www.lmfdb.org/EllipticCurve/Q/234446/a/1}{234446.a1}.
\end{thebibliography}
\end{document}

